\documentclass[11pt]{article}
\usepackage[margin=0.85in]{geometry}
\usepackage[T1]{fontenc}
\usepackage{lmodern}
\usepackage{amsmath,amssymb,booktabs,longtable,array,xurl,hyperref}
\hypersetup{colorlinks=true,linkcolor=blue,urlcolor=blue}
\setlength{\emergencystretch}{3em}
\newcommand{\dd}{\mathrm d}
\newcommand{\cp}{\mathord{>}p}
\title{Connecting the Existing LIFEORIG Prebiotic Network to an RNA Network}
\author{Concrete species handoff, reaction ownership, and coupled equations}
\date{November 1, 2026 --- revised coupling specification}
\begin{document}
\raggedbottom
\maketitle

\section{Direct answer: yes, with a chemically typed interface}
\textbf{Your existing network already contains a route from small precursors to
an activated pyrimidine ribonucleotide pool. Use that computed output to feed a
second network. Do not restart the RNA calculation with arbitrary monomer
concentrations.} The immediate, chemically closest second network is cyclic
nucleotide oligomerization. A template-copying network requires additional,
explicit substrate chemistry and complementary bases.

The audited source is
\path{reference_reactions/EARTH_LIKE/earth_like_volcanic_rock_combined_network_v1.txt},
which includes the prebiotic literature and ChemOrigins supplements. The
component file
\path{reference_reactions/PREBIOTIC/prebiotic_cyanosulfidic_literature_network_v1.txt}
contains the same named PBV1 pathway. The standalone
\path{prebiotic_polymer_autocatalytic_network_v2.txt} is a different hypothesis
scaffold: its HCN/formose/peptide rules do not substitute for the PBV1 nucleotide
pathway. Which file is loaded therefore matters.

The concrete connection is
\begin{equation}
\boxed{\text{existing PBV1-027--030}\ \longrightarrow\
\text{typed cyclic C/U nucleotides}\ \longrightarrow\
\text{new oligomerization network}.}
\end{equation}
For conventional 2-aminoimidazole-mediated copying, the connection instead is
\begin{equation}
\text{existing outputs}\ \longrightarrow\
\underbrace{\text{5'-NMP supply and activation}}_{\text{missing explicit bridge}}
\ \longrightarrow\ \text{template copying}.
\end{equation}
\textbf{The generic word ``activated'' does not make these two substrate types
interchangeable.} The following sections specify exactly where to cut the
networks, what to transfer, and what remains to implement.

\tableofcontents
\clearpage
\section{What your files already supply: exact reaction audit}
Use the following shorthand only in equations; retain canonical chemical IDs in
code: $G_2$ = glycolaldehyde, $G_3$ = glyceraldehyde, $Y$ = cyanamide,
$A_O$ = 2-aminooxazole, $P_A$ = pentose aminooxazoline,
$X$ = cyanoacetylene, $A_N$ = anhydronucleoside, and
$P_{\rm pyr}$ = the current pooled activated pyrimidine product.

\begin{longtable}{p{0.19\textwidth}p{0.43\textwidth}p{0.29\textwidth}}
\toprule
Existing ID & Existing equation, abbreviated & Meaning for the handoff \\
\midrule
PBV1-002,004 & glycolonitrile $\to G_2$ through the stored hydrolytic/photoredox
channels & Candidate upstream $G_2$ production; retain actual co-reactants. \\
PBV1-005,006 & $G_2+\mathrm{HCN}\to$ glyceronitrile; photoredox step $\to G_3$
& Candidate $G_3$ source. \\
VR166 & urea $\to Y+\mathrm{H_2O}$ & Cyanamide source candidate under its dry/hot conditions. \\
pbr-000743 & Ammonia + hydrogen cyanide $\to Y$ & Imported alternative cyanamide route;
its conditions and omitted chemistry need auditing. \\
pbr-000078 & acetylene + hydrogen cyanide $\to X$ & Imported cyanoacetylene source candidate. \\
PBV1-027 & $G_2+Y\to A_O$ & Entry to nucleotide synthesis, already present. \\
PBV1-028 & $A_O+G_3\to P_A$ & Sugar/base intermediate, already present. \\
PBV1-029 & $P_A+X\to A_N$ & Nucleoside intermediate, already present. \\
PBV1-030 & $A_N+\mathrm{phosphate}\to P_{\rm pyr}$ & Most downstream simple handoff in PBV1. \\
pbr-000708 & Orthophosphate + pbm-000227 $\to$ named cyclic cytidine product + water
& More specific candidate product identity than PBV1-030. \\
pbr-000048 & named cyclic cytidine product + water $\to$ named cyclic uridine product
+ ammonia & A C-to-U product branch is already represented in the combined file. \\
PBV1-026 & $5\,\mathrm{HCN}\to$ adenine & Supplies a base, not AMP or activated A. \\
\bottomrule
\end{longtable}
These are records in the files, not demonstrations that a simulation currently
produces nonzero yields. Their symbolic rate templates are not calibrated rate
constants. The accompanying \path{existing_reaction_evidence.tsv} records
29 relevant exact equations, source paths, line numbers, controls, and confidence
labels so this audit is checkable.

\subsection{Identity and duplication problems to resolve at this boundary}
\begin{enumerate}
\item \textbf{PBV1 and ChemOrigins overlap.} PBV1-027 and pbr-000043 describe the
same pathway step at different naming/detail levels; PBV1-028 overlaps
pbr-000044. Treat these as alternative representations until chemical identity
and mechanism show otherwise. Adding both as independent production channels
can double the apparent nucleotide supply.
\item \textbf{The cyclic endpoint is partly resolved already.} The named product
in pbr-000708 has a cyclic-phosphate structure. Its stereochemistry and linkage
identity must be checked against the source before assigning the canonical
$\beta$-D-ribocytidine 2',3'-cyclic phosphate ID. pbr-000048 supplies a corresponding
uridine branch. Do not claim this branch is absent.
\item \textbf{The intermediate chain is not automatically connected by names.}
pbr-000707 produces pbm-000559 whereas pbr-000708 consumes pbm-000227. The strings
are different. Their structures or an intervening reaction must establish the
connection; do not insert an identity conversion just to close the graph.
\item \textbf{Pools are not concrete molecules.} \texttt{pentose\_aminooxazoline}
can include isomers; \texttt{activated\_pyrimidine\_ribonucleotide} does not specify
base, phosphate position, anomer, or stereochemistry.
\item \textbf{Distinct compounds must stay distinct.} 2-aminooxazole is not
2-aminoimidazole; adenine is not adenosine or AMP; a 2',3'-cyclic phosphate is not
a 5'-phosphoroimidazolide. Generic \texttt{nucleotide}, \texttt{activated\_pyrimidine},
and \texttt{ribonucleotide\_precursor\_pool} are not automatic aliases of these.
\end{enumerate}

\section{Choose the handoff: use one of two cuts, not both}
\textbf{Recommended first cut: after PBV1-030.} Keep PBV1-027--030 in the existing
prebiotic network. Transfer its typed nucleotide products to the RNA network,
which owns oligomerization and subsequent RNA reactions. This directly tests
whether the network you already built supplies enough RNA material.

\textbf{Alternative cut: before PBV1-027.} Transfer the vector
\begin{equation}
\mathbf c_{\rm feed}=(C_{G_2},C_Y,C_{G_3},C_X,C_{\rm phosphate})
\end{equation}
to a second reactor which owns PBV1-027--030 plus the RNA chemistry. This is
appropriate if nucleotide synthesis requires different conditions from upstream
feedstock production. Remove those reactions from the upstream reactor in this
experiment. The second network is then ``nucleotide synthesis + RNA'', not just
RNA copying.

\subsection{Required handoff record}
At each output time or transfer event, export:
\begin{equation}
(t,\ \mathrm{species\_id},\ n_i\ [\mathrm{mol}],\ V\ [\mathrm L],\
T,\mathrm{pH},I,a_w,\mathrm{phase},\mathrm{identity\_status}).
\end{equation}
Also record relevant free-ion concentrations, radiation forcing, transfer
fraction/recovery, and reaction provenance. Concentration alone is insufficient
when reactor volumes differ. Export amounts remaining after all upstream
reactions and losses, not cumulative gross product formation.

For the PBV1 pool, use one of these explicit policies:
\begin{enumerate}
\item \textbf{Resolved chemistry:} replace PBV1-030 by specific C/U production
channels constrained by its source, with the UV/deamination branch and losses
explicit. The output has actual chemical IDs.
\item \textbf{Exploratory supply envelope:} retain the pool and introduce bounded
allocation fractions $f_C,f_U,f_{\rm other}$ with sum one. Then
$n_{C\cp}=f_C n_{P_{\rm pyr}}$ and $n_{U\cp}=f_U n_{P_{\rm pyr}}$ only as a
hypothesis about pool composition, not as an established chemical reaction.
Vary the fractions; do not report the result as a chemically resolved yield.
\end{enumerate}
If identity is unresolved and no allocation assumption is authorized in the
experiment configuration, the importer must reject direct copying-substrate
mapping rather than silently producing all four monomers.

\section{Coupling equations: no invented supply and no double counting}
Use mol and L for local amounts/volumes, seconds for time, and M for concentration.
Let $\mathbf n_P$ be upstream amounts, $\mathbf n_R$ RNA-reactor amounts,
and $\mathbf v_P,\mathbf v_R$ rate densities in M/s.

\subsection{Same pond: one chemical state, two reaction blocks}
If both networks describe the same liquid, combine their reaction blocks in one
state vector. A shared species has one amount and one ID:
\begin{equation}
\dot{\mathbf n}=V(S_P\mathbf v_P+S_R\mathbf v_R)
+\mathbf F_{\rm ext}-\mathbf F_{\rm loss}.
\label{same}
\end{equation}
For example, with $v_{30}$ the specific nucleotide-source rate and $u_C$ its
consumption by the RNA block,
\begin{equation}
\dot n_{C\cp}=V(v_{30,C}-u_C-v_{\rm hydro,C})+F_C^{\rm ext}.
\end{equation}
No artificial transfer coefficient is needed. The RNA block consumes the same
nucleotide inventory that the prebiotic block produces. This is the simplest
first implementation if the environmental conditions are compatible.

\subsection{Different reactors: conservative transfer}
For a well-mixed effluent flow $q$ in L/s from the prebiotic reactor,
\begin{align}
\dot{\mathbf n}_P&=V_PS_P\mathbf v_P+\mathbf F_P^{\rm in}
-q\mathbf C_P-\mathbf F_P^{\rm other\ loss},\\
\dot{\mathbf n}_R&=V_RS_R\mathbf v_R+\eta qH\mathbf C_P
-q_R\mathbf C_R-\mathbf F_R^{\rm other\ loss}.
\label{transfer}
\end{align}
$H$ maps chemically identical transferred species to receiving IDs;
$\eta\in[0,1]$ is a specified recovery for a simple common-recovery model. Track
$(1-\eta)qH\mathbf C_P$ as transfer loss and other effluent species in the
receiving mixture or a discarded-stream ledger. Selective recovery needs a
physical separation model. If amounts accumulate, evolve volumes too:
$\dot V_P=q_{P,\rm in}-q$ and $\dot V_R=q-q_R-E_R$ (including other solvent terms).
Do not hold volume fixed without the corresponding compensating flows.

For a batch transfer of fraction $\alpha$ at $t_m$:
\begin{align}
\mathbf n_P^+&=(1-\alpha)\mathbf n_P^-,\\
\mathbf n_R^+&=\mathbf n_R^-+\eta H\alpha\mathbf n_P^-,\\
C_{i,R}^+&=n_{i,R}^+/V_R^+.
\end{align}
An unseparated liquid aliquot also transfers $\alpha V_P$ solvent. Concentrating
it into a smaller volume requires an explicit evaporation/separation step.

\subsection{Offline replay is only a preliminary coupling test}
A stored upstream trajectory can drive a downstream calculation. It is a
one-way reservoir approximation unless the exported material was actually
removed in the upstream calculation. Do not feed a downstream sink with
$\max(\dot n_i,0)$: net accumulation is not a delivery process. For final coupled
results use Eq.~\eqref{same} or Eq.~\eqref{transfer}, or conservative batch events.

\clearpage
\section{Second network A: directly use the cyclic pyrimidine outputs}
\textbf{This is the nearest concrete extension of your present prebiotic network.}
Define $M_C=C\cp$ and $M_U=U\cp$ as the verified 2',3'-cyclic ribonucleotides.
Cyclic-nucleotide wet--dry oligomerization is an experimentally studied branch
\cite{dass,cyclic}. It does not require pretending these species are activated
5'-NMPs. A C/U-only calculation tests oligomer formation; it does not supply a
closed Watson--Crick complementary copying alphabet.

\subsection{Concrete additional channels}
Let $R_s^{\cp}$ be an oligomer with an explicitly tracked terminal cyclic
phosphate, a free incoming 5'-OH partner, and stored internal linkage identities.
For $a,b\in\{C,U\}$ add the following effective reaction rules:
\begin{longtable}{p{0.16\textwidth}p{0.48\textwidth}p{0.27\textwidth}}
\toprule
New ID & Effective channel & Parameter/data required \\
\midrule
RNA-CYC-01 & $M_a+M_b\to R_{ab}^{\cp}$ & Dimerization, sequence and linkage branching. \\
RNA-CYC-02 & $R_s^{\cp}+M_b\to R_{sb}^{\cp}$ & End-specific cyclic-phosphate ring-opening addition. \\
RNA-CYC-03 & $M_b+\mathrm{H_2O}\to b\text{-}2'\mathrm{MP}$ & Competing ring opening branch. \\
RNA-CYC-04 & $M_b+\mathrm{H_2O}\to b\text{-}3'\mathrm{MP}$ & Competing ring opening branch. \\
RNA-CYC-05 & $R_s\to$ explicit fragments & Linkage-dependent cleavage and fragment ends. \\
RNA-CYC-06 & $R_s+\mathrm{site}\rightleftharpoons B_s$ & Optional retention/desorption and site budget. \\
\bottomrule
\end{longtable}
These are \emph{new proposed coarse-grained rules}, not claimed elementary
mechanisms already implemented. Resolve chain orientation and termini before
compiling them. Ring-opening transesterification does not generically release
one water per new bond: do not reuse the existing ``nucleotide + nucleotide
$\to$ dinucleotide + water'' condensation proxy for this mechanism.

For a deliberately reduced length model, $O_n$ denotes the concentration of
chains of length $n$ with an extendable cyclic end. Set $O_1=M$ for a single
monomer type. Example event-rate densities are
\begin{align}
v_1&=k_1O_1^2,&v_n&=k_nO_nO_1\quad(n\ge2),\\
\dot O_1\big|_{\rm growth}&=-2v_1-\sum_{n\ge2}v_n,\\
\dot O_2\big|_{\rm growth}&=v_1-v_2, &
\dot O_n\big|_{\rm growth}&=v_{n-1}-v_n\quad(n\ge3).
\end{align}
Concentration dynamics also include $-O_n\dot V/V$, sources, cleavage, and
inactivation. $k_n$ is M$^{-1}$s$^{-1}$ in the aqueous effective model. Dry-film
kinetics need a calibrated phase-specific treatment rather than divergent
$1/V$ concentrations at zero volume. Mixed sequences require separate $O_s$
and separate 2'-5'/3'-5' linkage branches; a length-only chain is not a ribozyme.
At the maximum stored length, either stop extension consistently or send mass
to a tracked longer-chain bin.

\subsection{The first result to compute}
Run the prebiotic network and this block together. Start with zero oligomers.
Measure, in nucleotide-equivalent mol,
\begin{equation}
B_{\ge L}(t)=\sum_{|s|\ge L}|s|n_{R_s}(t),\qquad
Y_{\ge L}(t)=\frac{B_{\ge L}(t)}{B_{\rm delivered}(t)}.
\end{equation}
Count delivered material once and include initial inventory in the denominator
if nonzero. Report $L=2,5,10$ as diagnostic cutoffs, not functional thresholds.
The test is whether the actual upstream supply produces an oligomer inventory
that survives subsequent wet phases. Compare with a control that suppresses
PBV1-030 production: the downstream material should then disappear or exhaust
its stated initial stock. This directly answers the precursor-feeding question
before adding speculative replication.

\section{Second network B: an explicit template-copying branch}
\subsection{Required inlet species and the current gap}
For the well-studied 2-aminoimidazole route, the inlet is chemically resolved
5'-NMPs and an activation system, or already activated 5'-nucleotides. Define
$N_b$ = 5'-NMP, $A_b$ = its 2AI-activated derivative, $I$ = 2AI, and
$B_{bc}$ = a 5'-5' imidazolium-bridged dinucleotide. These are different from
$C\cp,U\cp$. The copying mechanism involves the bridged intermediate
\cite{walton,activation}.

The audited PBV1 endpoint does \textbf{not} establish a conversion
$C\cp\to C\text{-}5'\mathrm{MP}\to A_C$. Simple opening of a 2',3'-cyclic phosphate
gives 2'/3'-phosphates, not a 5'-phosphate. A conversion module would require
supported positional/activation chemistry with its own material and reagent
budget. Alternatively choose a separately supported 5'-NMP synthesis branch.
Neither conversion is supplied by renaming the pool.

Likewise, the PBV1 adenine product is not an A-nucleotide supply. For a
C-containing template, copying needs G; for a U-containing template, copying
needs A. A C/U-only product stream cannot copy itself through canonical pairing.
Find and validate A/G ribonucleotide routes in the larger imported mechanism,
or mark A/G feeds as external in a conditional experiment. DNA products such
as deoxyadenosine cannot silently fill this RNA gap.

\subsection{Minimal additional reaction block}
A reduced copying block, after a validated 5'-substrate inlet exists, is
\begin{align}
N_b+\mathrm{activation\ reagents}&\longrightarrow A_b+\mathrm{spent\ products},\\
A_b+A_c&\rightleftharpoons B_{bc}+I,\\
A_b+\mathrm{H_2O}&\longrightarrow N_b+I,\\
X_{T,u}+B_{bc}&\rightleftharpoons Y_{T,u,bc},\\
Y_{T,u,bc}&\longrightarrow X_{T,ub}+A_c,\\
X_{T,\bar T}&\longrightarrow D_T,\\
D_T&\rightleftharpoons T+\bar T.
\end{align}
$X_{T,u}$ is a template--primer complex and $Y$ its substrate-bound complex.
The bridge-mediated step adds one base and releases an activated nucleotide;
it does not incorporate both bases as a covalent RNA dinucleotide. The first
line is a module requirement, not a balanced invented activation reaction.
Use the selected experimental mechanism for its actual reagents and rates.
Add mismatches, stalled complexes, hydrolysis, and primer binding explicitly.

A first benchmark can copy a seeded G-rich template using a typed activated-C
feed. It tests one copying leg. It does not demonstrate C/G self-replication
unless the reverse leg has a G-substrate supply, initiation, strand release,
and sustained turnover. This is a useful conditional benchmark \emph{after}
the inlet chemistry is specified, not a direct consequence of PBV1-030.

\section{Calculate whether the upstream supply is sufficient}
\subsection{Inventory bound before any kinetic simulation}
For an RNA sequence $s$ of length $L$, let $\ell_b(s)$ be its number of base $b$
and $A_b$ the mol of usable substrate supplied after transfer and activation
losses. The maximum mol of newly synthesized copies is bounded by
\begin{equation}
n_{s,\max}\le\min_{b:\ell_b>0}\frac{A_b}{\ell_b(s)}.
\label{bound}
\end{equation}
Use the composition of the \emph{new complementary strand}, not the template,
for one-leg copying; use the combined two-strand composition for a full pair.
If a required substrate has zero supply, the bound is zero regardless of the
chosen copying rate. This immediately exposes missing A/G feeds.

\subsection{Worked transfer example: illustrative numbers, not a prediction}
Suppose a completed upstream run reports 10 $\mu$M typed cyclic pyrimidines in
10 mL, or 100 nmol total. Transfer the full batch with recovery 0.4, then
explicitly concentrate it to 1 mL. The receiving amount is 40 nmol and the
concentration is 40 $\mu$M. With an \emph{assumed} 60:40 C/U composition:
\begin{equation}
n_{C\cp}=24\ \mathrm{nmol},\quad n_{U\cp}=16\ \mathrm{nmol}.
\end{equation}
This feeds network A. Even if all of it became 10-mers, the absolute upper
bound is $40/10=4$ nmol chains; actual yields are lower and sequence-dependent.
It cannot directly feed network B. Under an additional hypothetical conversion
to usable C/U copying substrates, copying a CU-containing template still
requires external A/G; without them Eq.~\eqref{bound} remains zero for that leg.
Every number in this example is arithmetic for checking the coupling, not a
literature rate or an output already obtained from LIFEORIG.

\subsection{Flux bound for sustained reproduction}
For delivered usable substrate $F_b$ in mol/s, a target population $n_R$ mol of
strands, average new-strand base count $\bar\ell_b$, turnover $g$ s$^{-1}$,
and substrate loss flux $L_b$ mol/s:
\begin{equation}
F_b\ \ge\ \bar\ell_b g n_R+L_b
\end{equation}
is a necessary resource condition at approximately steady inventory. It is not
sufficient: binding, mismatches, strand release, and competing reactions can
still block replication. If conversion to the copying substrate has recovery
$\chi_b$, use $F_b=\eta_b\chi_bF_b^{\rm upstream}$ and report how $\chi_b$ is
obtained. An unknown bridge yield is a sensitivity parameter, not proof of
chemical connectivity.

\section{Exact implementation procedure for this repository}
\begin{enumerate}
\item \textbf{Freeze the upstream file and audit identities.} Use the combined
Earth-like file as the initial scenario. Resolve aliases such as
Glycolaldehyde/glycolaldehyde and Phosphate/orthophosphate only with verified
phase/speciation semantics. Remove overlapping pathway representations from
the selected kinetic model; retain their provenance.
\item \textbf{Obtain a real upstream trajectory.} Ensure the selected prebiotic
channels actually compile and have defined rate laws. Save amounts and volumes
for the boundary species, with all competing sinks active. Presence in a text
file does not establish solver support or production.
\item \textbf{Type the endpoint.} Refine PBV1-030 into the verified cyclic C/U
products, using pbr-000708/048 and primary chemistry as evidence. Until resolved,
run only explicitly labeled pool-allocation sensitivity studies.
\item \textbf{Add network A first.} Implement RNA-CYC-01--05 with measured or
explicitly explored effective coefficients, chemically correct termini, and
tracked linkage classes. Initialize oligomers at zero. Keep precursor production
in the existing network; do not seed a second arbitrary monomer inventory.
\item \textbf{Couple using one state or conservative flows.} Same pond:
Eq.~\eqref{same}. Different conditions/reactors: Eq.~\eqref{transfer} or batch events.
Transfer salts and inhibitors too unless a separation step is explicitly modeled.
\item \textbf{Produce the first decision plot.} For actual upstream histories,
plot delivered nucleotide amount, cyclic monomer loss, $B_{\ge2},B_{\ge5},B_{\ge10}$,
and material residual versus time. Repeat over uncertain production and
polymerization coefficients. This answers whether your prebiotic output can
sustain an RNA-oligomer population.
\item \textbf{Only then select the copying chemistry.} Supply a validated
substrate bridge and complementary alphabet, or label their feeds as external.
Add the bridge-intermediate and template-complex reactions above. Separate
seeded primer extension from a complete reproduction cycle.
\end{enumerate}

\subsection{Specific code touchpoints}
\begin{itemize}
\item \path{src/chemical_types/molecule_model.py}: a \texttt{ReferenceMolecule}
has sequence, phase, state, and conformation fields. Extend the identity scheme
to distinguish base, sugar stereochemistry, phosphate position, activation group,
and oligomer termini; generic aliases must not merge these states.
\item \path{src/reactions/reaction_rate_models.py}: implement aqueous/film,
cyclic-growth, and copying rate models explicitly. Existing symbolic labels
such as \texttt{chemorigins\_experimental} are not numeric kinetics by themselves.
\item \path{src/chem_network/chem_system_builder.py}: report which audited
channels compile. It currently collects unsupported reactions; a coupling run
must fail if its required supply or RNA channels are missing.
\item \path{src/chem_network/chemical_network.py}: register shared species once
and add the upstream and RNA reaction blocks. Use one stoichiometric inventory
for a same-pond experiment.
\item \path{src/environment/surf_pond/surface_pond_env.py}: use the pond volume,
forcing, and transport state. Its development exit must be resolved before
claiming an end-to-end pond calculation.
\end{itemize}
The output of this work is a concrete coupling specification and source audit;
the new RNA rules and their kinetic parameterization are not implemented by
this document.

\section{Acceptance checks and the claim each permits}
\begin{longtable}{p{0.48\textwidth}p{0.43\textwidth}}
\toprule
Required check & Permitted conclusion \\
\midrule
All required source reactions compile; typed cyclic products have nonzero
amounts with budget closure & The selected prebiotic model supplies nucleotide
material under the stated parameters. \\
RNA-CYC channels convert that supply into tracked oligomers, with linkage and
loss accounting & Prebiotic products feed an RNA-oligomer network. \\
Template extension consumes a validated copying-substrate stream with sequence
and mismatch tracking & Conditional template-directed copying occurs. \\
Both complementary legs, strand release, and replenishment maintain descendants
across cycles & Conditional sustained RNA reproduction occurs. \\
Inherited sequence changes cause differential reproduction under matched
conditions across replicated lineages & The model exhibits RNA sequence selection. \\
\bottomrule
\end{longtable}
For closed internal chemistry, check $BS=0$ using an elemental matrix $B$.
For the RNA block also track nucleotide equivalents in monomers, intermediates,
chains, and fragments. Sources and discarded material must close the full
budget. Reducing volume must increase concentration without creating amount.
No-precursor and missing-required-base controls must respect Eq.~\eqref{bound}.
Do not interpret a large generic polymer count as a functional ribozyme.

\textbf{Decision for the present repository:} yes, connect the existing products
to a second network. Start at the PBV1-030 cyclic pyrimidine endpoint, resolve
its identity using the more detailed records, and implement cyclic-nucleotide
oligomerization. Direct input into a conventional activated-5'-nucleotide copying
network is currently an unfilled chemical bridge, not a mere software adapter.

\begin{thebibliography}{9}
\bibitem{powner} M. W. Powner, B. Gerland, and J. D. Sutherland.
Synthesis of activated pyrimidine ribonucleotides in prebiotically plausible
conditions. \emph{Nature} 459, 239--242 (2009).
\url{https://doi.org/10.1038/nature08013}.
\bibitem{patel} B. H. Patel et al. Common origins of RNA, protein and lipid
precursors in a cyanosulfidic protometabolism. \emph{Nature Chemistry} 7,
301--307 (2015). \url{https://doi.org/10.1038/nchem.2202}.
\bibitem{dass} RNA Oligomerisation without Added Catalyst from 2',3'-Cyclic
Nucleotides by Drying at Air--Water Interfaces. \emph{ChemSystemsChem} (2023).
\url{https://doi.org/10.1002/syst.202200026}.
\bibitem{cyclic} High-Yield Prebiotic Polymerization of 2',3'-Cyclic Nucleotides
under Wet--Dry Cycling. \emph{ACS Central Science} (2025).
\url{https://doi.org/10.1021/acscentsci.5c00488}.
\bibitem{walton} T. Walton and J. W. Szostak. A Kinetic Model of Nonenzymatic
RNA Polymerization by Cytidine-5'-phosphoro-2-aminoimidazolide.
\emph{Biochemistry} 56 (2017).
\url{https://doi.org/10.1021/acs.biochem.7b00792}.
\bibitem{activation} Potentially Prebiotic Activation Chemistry Compatible
with Nonenzymatic RNA Copying (2020).
\url{https://pmc.ncbi.nlm.nih.gov/articles/PMC9594304/}.
\end{thebibliography}
\end{document}
